\documentclass[conference]{IEEEtran}
\IEEEoverridecommandlockouts
% The preceding line is only needed to identify funding in the first footnote. If that is unneeded, please comment it out.
%Template version as of 6/27/2024

\usepackage{cite}
\usepackage{svg}
\usepackage{amsmath,amssymb,amsfonts}
\usepackage{algorithmic}
\usepackage{graphicx}
\usepackage{textcomp}
\usepackage{xcolor}
\def\BibTeX{{\rm B\kern-.05em{\sc i\kern-.025em b}\kern-.08em
    T\kern-.1667em\lower.7ex\hbox{E}\kern-.125emX}}
\newcommand{\linebreakand}{%
  \end{@IEEEauthorhalign}
  \hfill\mbox{}\par
  \mbox{}\hfill\begin{@IEEEauthorhalign}
}
% ────────────────────────────────────────────────────────────────────────────
\begin{document}

% ── Title block ─────────────────────────────────────────────────────────────
\title{Aether-Edge: Edge-Native Predictive HVAC Control via Gossip Networks}

\author{%
  \IEEEauthorblockN{Kassandra Langlois}
  \IEEEauthorblockA{%
    MEng ECE. University of Waterloo\\k2langlois@uwaterloo.ca
  }
  \and
  \IEEEauthorblockN{Prabhjyot Singh}
  \IEEEauthorblockA{%
    MEng ECE. University of Waterloo\\p269singh@uwaterloo.ca
  }
  \and
  \IEEEauthorblockN{Rick Liu}
  \IEEEauthorblockA{%
    MEng ECE. University of Waterloo\\m385liu@uwaterloo.ca
  }
}

\maketitle
\thispagestyle{plain}
\pagestyle{plain}
% ── Abstract ────────────────────────────────────────────────────────────────
\begin{abstract}
Allocating a contentioned shared resource across competing nodes in a sensor
network is fundamentally limited by the freshness of each node's view of the
global demand landscape.  Centralized arbitration introduces an Age of
Information (AoI) gap — the controller acts on data that is already tens of
seconds old — causing suboptimal allocation during rapid transients.  We
present Aether-Edge, a decentralized edge-native framework in which each
sensor node independently estimates a Time-To-Impact (TTI) urgency signal
from a local rolling buffer via ordinary least-squares regression, and
propagates it through a threshold-triggered $n$-hop gossip protocol.  Each
actuator allocates its proportional share of the contentioned resource using
only the urgency landscape it has accumulated from neighborhood gossip, with
AoI\,=\,0.  We demonstrate the approach on building HVAC, where multiple
thermal zones compete for a shared cooling budget $Q_{\text{total}}$ set
below the sum of per-zone capacities to enforce genuine contention.
Benchmarked against a centralized polling baseline using the identical
urgency algorithm, Aether-Edge achieves \todo{X}\% reduction in cumulative
comfort violation, \todo{Y}\% change in energy consumption, and \todo{Z}\%
fewer gossip messages, while reducing mean AoI from \todo{W}\,s to zero.
\end{abstract}

% ============================================================================
\section{Introduction}
% ============================================================================

\subsection{Problem Statement}

Buildings account for approximately 40\% of global primary energy consumption,
with HVAC systems responsible for up to 50\% of total building energy use
\cite{iea2023}.  Despite this outsized share, the control strategies deployed
in the vast majority of commercial and institutional buildings remain reactive
and centralized.  Conventional Building Automation Systems (BAS) operate on a
polling architecture: temperature sensors report readings to a central
controller over a fieldbus or IP network, the controller runs a PID or
schedule-based algorithm, and actuator commands are pushed back to Variable
Air Volume (VAV) dampers --- all at intervals typically ranging from 5 to 30
seconds \cite{ashrae2017}.

This architecture has three fundamental limitations.  First, thermal comfort
violations occur because the system can only react to conditions that have
already developed: a lecture hall filling with students, a server room
absorbing a burst of compute load, or a conference room heated by afternoon
solar gain will breach its setpoint before the next poll cycle fires a
corrective command.  Second, the spatial heterogeneity of real building loads
--- different occupancy densities, equipment profiles, and solar exposures
across zones --- means that a single fixed cooling schedule wastes energy on
empty rooms while under-serving occupied ones.  Third, and most fundamentally,
the centralized controller makes every decision based on \emph{stale data}:
the environment has continued to evolve during the polling interval, network
transit, and compute delay, so the commanded actuation is mis-calibrated to
the state the building was in, not the state it is in.

\subsection{Motivation}

The key insight motivating Aether-Edge is that the information required to
make a good cooling decision --- specifically, how fast a zone is warming and
how close it is to its setpoint --- can be computed \emph{locally and
instantaneously} from the sensor's own rolling history.  No network round-trip
is needed.  A node that observes its zone temperature rising at
$\dot{T} = 0.05\,^\circ\text{C/s}$ with 2\,°C of headroom before its setpoint
can immediately compute that a comfort breach is approximately 40 seconds away.
This Time-To-Impact (TTI) estimate is a pure local computation with zero Age
of Information.

The Age of Information (AoI) framework, first formalized by Kaul~et~al.
\cite{kaul2012}, provides the right metric to quantify this staleness gap.
AoI measures the elapsed time since the most recently delivered measurement
was generated at the source.  In a BAS with a 15-second poll interval,
2-second network jitter, and 3-second compute delay, the mean AoI at actuation
time is approximately 19.4\,s --- nearly half a minute of environmental
evolution that the controller is blind to.  By contrast, an edge-native node
that infers from its current sensor reading has AoI\,=\,0.

Gossip protocols \cite{kempe2003, boyd2006} provide a lightweight mechanism
for zones to share their local TTI estimates with spatial neighbors, enabling
the building's shared cooling budget to be allocated proportionally to urgency
without a central coordinator.  Because gossip propagates only when a zone
detects a meaningful rate of change (threshold-triggered), the communication
overhead is far lower than unconditional periodic polling.

A concrete scenario motivates the design: a university building with two
classrooms, a computer lab, a conference room, and a faculty office follows
heterogeneous occupancy schedules.  A morning lecture in Classroom~101 ramps
heat load by 0.015\,°C/s while the faculty office remains near ambient.  A
centralized controller with a 15-second poll cycle will not issue a corrective
cooling command to Classroom~101 until the \emph{next} poll fires after the
load appears --- roughly 7.5 seconds late on average, and up to 15 seconds
late in the worst case.  The edge node detects the rising slope within one
buffer window and pre-emptively increases the damper opening before the
setpoint is breached.

\subsection{Goal and Objectives}

The broader thesis is that \emph{urgency-based $n$-hop gossip at the edge
is a general mechanism for allocating contentioned shared resources in
sensor networks}: each node independently infers how urgently it needs a
resource (via TTI), gossips that urgency to neighbors to form a local view
of the global demand, and allocates proportionally with zero data staleness.
The HVAC environment --- where $Q_{\text{total}}$ is the contentioned resource
and thermal zones are the competing nodes --- is a concrete demonstration
substrate for this principle.  The specific thesis instantiated here is that
a decentralized edge network outperforms a centralized reactive controller on
comfort, energy, and network overhead, with the gap attributable directly to
the AoI difference between the two architectures.

The specific measurable objectives are:
\begin{enumerate}[]
  \item Design and implement a 3D thermal diffusion physics engine that
        accurately models heat diffusion, occupancy loads, and VAV damper
        cooling with numerical stability enforced by the CFL condition.
  \item Design a per-zone inference engine using a rolling circular buffer
        and ordinary least-squares slope estimation to compute $dT/dt$ and
        TTI without a Kalman filter or external model.
  \item Implement a threshold-triggered urgency gossip protocol for
        decentralized proportional allocation of a shared cooling budget
        across zones.
  \item Implement a configurable centralized baseline with polling interval,
        network jitter, and compute delay to reproduce realistic BAS
        behaviour.
  \item Benchmark edge versus centralized control on five metrics: cumulative
        comfort violation, cumulative energy, max overshoot, mean AoI, and
        total message count.
\end{enumerate}

% ============================================================================
\section{Related Work}
% ============================================================================

\subsection{Building Automation and Centralized HVAC Control}

The dominant communication protocol for centralized BAS is ASHRAE Standard
135 (BACnet) \cite{ashrae2017}, which defines a polling-based message model
in which a central controller periodically reads object properties (such as
zone temperature) from field devices.  While BACnet supports event-driven
messaging, practical deployments overwhelmingly use polled data collection at
intervals of 5--30 seconds to limit network traffic.

Classical thermostat control uses PID feedback: the error between measured
temperature and setpoint drives a proportional-integral-derivative actuator
command.  PID is simple and robust for single-zone systems, but in multi-zone
buildings with a shared cooling plant it suffers from integral windup,
inter-zone coupling, and an inability to anticipate load changes before they
appear in the measurement.

Model Predictive Control (MPC) for HVAC was proposed by Oldewurtel~et~al.
\cite{oldewurtel2012} as a principled solution: a central controller solves a
rolling-horizon optimization problem that minimizes energy cost subject to
comfort constraints, using a calibrated thermal model and weather forecasts.
Ma~et~al.~\cite{ma2012} extend this to stochastic MPC, incorporating
uncertainty in occupancy and weather.  These approaches are effective but
impose a significant computational and communication burden: the central
optimizer requires a global model of all zones, continuous sensor polling to
update state estimates, and real-time solution of a quadratic program at each
timestep.  In large buildings, the polling and model-update overhead itself
contributes to AoI.  Aether-Edge sidesteps the global model requirement
entirely by replacing optimization with local TTI inference.

\subsection{Decentralized and Multi-Agent Control}

Multi-agent approaches distribute the HVAC control problem across autonomous
zone-level agents.  Hao~et~al.~\cite{hao2012} demonstrate that zone agents
negotiating via local price signals can achieve near-optimal demand-response
performance without a central optimizer.  Olfati-Saber~et~al.~\cite{olfati2007}
establish the theoretical foundation for consensus in networked multi-agent
systems: under mild connectivity conditions, gossip-based averaging converges
to a global consensus in $O(\log N)$ rounds for $N$ agents, with convergence
rate governed by the second-smallest eigenvalue of the graph Laplacian
(spectral gap).

Despite these advances, most multi-agent HVAC work still uses
\emph{synchronous} periodic communication --- agents exchange messages at
fixed intervals regardless of whether anything has changed.  This is
functionally equivalent to centralized polling from an AoI perspective.  The
key innovation in Aether-Edge is that gossip is \emph{event-triggered}: a
zone transmits only when its $|dT/dt|$ exceeds a configurable threshold,
directly coupling communication frequency to the rate at which conditions
change.  During stable periods (empty rooms, steady-state temperatures) the
network is silent; during transient events (occupancy surges, equipment
startups) gossip propagates urgency at the rate the environment actually
demands.

\subsection{Age of Information in Networked Control}

AoI was introduced by Kaul, Yates, and Gruteser \cite{kaul2012} as a metric
for the freshness of status updates from the perspective of the destination: it
measures the time elapsed since the generation of the most recently delivered
packet.  Sun~et~al.~\cite{sun2017} derive AoI-optimal update policies for
single-source systems, showing that the optimal transmission rate balances
freshness against network congestion.  Kaul and Yates \cite{kaul2020} extend
this to energy-constrained sensors, establishing that threshold-based
transmission policies (transmit only when the change in state exceeds a
threshold) minimize AoI per unit energy, which directly justifies our
\texttt{talk\_threshold} design.

In cyber-physical control systems, high AoI translates directly to control
degradation.  Costa and Ephremides \cite{costa2016} demonstrate analytically
that stale control actions in discrete-time linear systems cause setpoint
overshoot proportional to the AoI and the open-loop dynamics gain.  Applied
to HVAC, this means a controller acting on 19-second-old data in a zone
warming at $0.05\,^\circ\text{C/s}$ effectively over-estimates the zone's
current temperature by $19 \times 0.05 \approx 1\,^\circ\text{C}$ --- enough
to flip a comfort decision.  Aether-Edge eliminates this error by collapsing
AoI to zero.

\subsection{Gossip Protocols and Sensor Networks}

Kempe, Dobra, and Gehrke \cite{kempe2003} established gossip as a general
mechanism for distributed aggregate computation (sum, average, percentile)
in sensor networks, with convergence guarantees for connected graphs.  Boyd
et~al.~\cite{boyd2006} analyzed randomized gossip and derived tight bounds on
mixing time as a function of graph topology: for a $k$-regular graph on $N$
nodes, average consensus requires $O(N \log N / k)$ rounds.  For the
university floor topology in our experiments ($N \approx 40$, $k \approx 3$),
two gossip rounds provide sufficient urgency propagation to all spatial
neighbors.

Threshold-based gossip --- where a node transmits only when its state changes
by more than a threshold $\tau$ --- was analyzed in the context of event-driven
sensor networks by Rabbat and Nowak \cite{rabbat2004}: the total number of
transmissions scales as $O(N/\tau^2)$ rather than $O(N \log N)$ for full
consensus, providing a direct bandwidth-accuracy tradeoff.  Our
\texttt{talk\_threshold} parameter implements exactly this tradeoff: setting
it to $0.01\,^\circ\text{C/s}$ suppresses gossip during all stable periods
while ensuring that meaningful thermal trends are propagated within two rounds.

% ============================================================================
\section{Methodology}
% ============================================================================

\subsection{System Architecture}

Aether-Edge is organized as seven functional blocks with a strict information
flow hierarchy:

\begin{figure}[h]
  \centering
  \includegraphics[width=1.00\linewidth]{system.jpeg}
  \caption{Aether-Edge block architecture.}
  \label{fig:block-arch}
\end{figure}

\textbf{Block~1 (World)} is the 3D thermal physics engine.  It maintains the
temperature field $T \in \mathbb{R}^{N_x \times N_y \times N_z}$ and advances
it by one timestep on each call to \texttt{world.step()}.  \textbf{Block~2
(Visualization)} renders temperature heatmaps and damper states as animated
GIFs for qualitative analysis.  \textbf{Block~3 (Interface)} is the
\emph{sole I/O boundary}: it reads zone temperatures from the World and writes
damper opening fractions back to it.  It is the only block that touches the
World directly, enforcing complete decoupling between the physics engine and
the inference network.  \textbf{Block~4 (SensorNetwork)} maintains the mesh
communication graph and computes topology metrics.  \textbf{Block~5
(SensorField)} orchestrates per-node rolling buffer updates, TTI inference,
and gossip rounds.  \textbf{Block~6 (Simulation)} is the main loop that
sequences World and Interface calls for both edge and centralized modes.
\textbf{Block~7 (Benchmark)} runs both modes from identical initial conditions
and computes the comparison metrics.

The key design principle is that the SensorField block \emph{never} receives
the World's temperature array.  The Interface extracts scalar
$(t_k,\, T_k)$ pairs at each sensor position and passes them to the
SensorField as opaque scalar streams.  This enforces domain-agnostic inference:
replacing the temperature interface with one that reads CO$_2$ or humidity
requires no changes to Blocks~4--7.

\subsection{Thermal Physics Engine}

The thermal environment is governed by the heat diffusion PDE:

\begin{equation}
  \frac{\partial T}{\partial t} =
      \alpha \nabla^2 T
      + Q_{\text{heat}}
      - Q_{\text{cool}},
  \label{eq:pde}
\end{equation}

where $\alpha = 0.02\,\text{m}^2/\text{s}$ is the thermal diffusivity of air
and $Q_{\text{heat}}$, $Q_{\text{cool}}$ are spatially distributed heat source
and sink terms, respectively.  The Laplacian is discretized using a
second-order central-difference scheme on a 3D Cartesian grid with
$\Delta x = 1\,\text{m}$ resolution and zero-flux (Neumann) boundary
conditions at all domain walls.

The timestep is selected automatically to satisfy the 3D FTCS diffusion
stability criterion with a safety factor of 0.4:

\begin{equation}
  \Delta t = 0.4 \times \frac{\Delta x^2}{6\alpha}.
  \label{eq:dt}
\end{equation}

Heat sources $Q_{\text{heat}}$ are specified per zone with configurable
occupancy and equipment schedules.  Each zone $i$ has a VAV damper with a per-zone maximum flow capacity
$f_i^{\max}$.  The damper opening fraction $A_i \in [0,1]$ determines the
requested flow $q_i = A_i \cdot f_i^{\max}$.  Because the shared plant
capacity $Q_{\text{total}}$ is set below the sum of all per-zone capacities
($Q_{\text{total}} < \sum_i f_i^{\max}$), contention is guaranteed whenever
multiple zones are simultaneously active.  The actual delivered cooling is:

\begin{equation}
  Q_{\text{cool},i} = A_i \cdot f_i^{\max} \cdot
    \min\!\left(1,\; \frac{Q_{\text{total}}}{\sum_j A_j \cdot f_j^{\max}}\right),
  \label{eq:cooling}
\end{equation}

so the total delivered cooling never exceeds $Q_{\text{total}}$, and zones
that open their dampers more aggressively receive a proportionally larger
share of the scarce budget.

Table~\ref{tab:zones} summarizes the five university floor zones used in
all experiments.

\begin{table}[t]
\centering
\caption{University Floor Zone Parameters}
\label{tab:zones}
\begin{tabular}{lccl}
\toprule
\textbf{Zone} & \textbf{Setpoint} & \textbf{Peak $Q_{\text{heat}}$}
              & \textbf{Schedule} \\
\midrule
Classroom 101    & 22\,°C & 0.015 & Morning lecture \\
Classroom 102    & 22\,°C & 0.014 & Afternoon lecture \\
Computer Lab     & 20\,°C & 0.018 & Always-on + occupancy \\
Conference Room  & 22\,°C & 0.012 & Scheduled meetings \\
Faculty Office   & 22\,°C & 0.002 & Light occupancy \\
\bottomrule
\end{tabular}
\end{table}

\subsection{Per-Zone Inference: Rolling Buffer and TTI}

Each sensor node maintains a circular rolling buffer of the most recent
$N = \lfloor t_{\text{buf}} / \Delta t \rfloor$ scalar readings, where
$t_{\text{buf}} = 30\,\text{s}$ is the buffer window.  Readings older than
$t_{\text{buf}}$ are discarded automatically as new measurements arrive.

The rate of change $\hat{\dot{T}}$ is estimated by an ordinary least-squares
(OLS) linear fit over the buffer contents:

\begin{equation}
  \hat{\dot{T}} = \frac{n \sum t_k T_k - \sum t_k \sum T_k}
                       {n \sum t_k^2 - \bigl(\sum t_k\bigr)^2},
  \label{eq:ols}
\end{equation}

where the sums run over the $n$ samples in the buffer.  OLS over the buffer is
preferred to a simple finite difference because it is robust to individual
noisy readings and to sensor dropout events within the window: a single
anomalous sample shifts the slope estimate by at most $1/n$ rather than
dominating the result.

The Time-To-Impact is then:

\begin{equation}
  \text{TTI} = \begin{cases}
    \dfrac{T_{\text{sp}} - T}{\hat{\dot{T}}} & \text{if } \hat{\dot{T}} > 0, \\[6pt]
    \infty & \text{otherwise,}
  \end{cases}
  \label{eq:tti}
\end{equation}

where $T_{\text{sp}}$ is the zone comfort setpoint.  TTI is finite only when
the zone is actively warming toward its setpoint.  The urgency signal is:

\begin{equation}
  u = \frac{1}{\text{TTI}},
  \label{eq:urgency}
\end{equation}

so that $u \to \infty$ as a breach becomes imminent and $u = 0$ when the zone
is stable or cooling.  Critically, both $\hat{\dot{T}}$ and $u$ are computed
entirely from the node's own current buffer --- the AoI of this signal is
identically zero.
% \subsection{Urgency Gossip and Resource Negotiation}

\subsection{Gossip Trigger Condition}

After computing its local urgency score, each sensor node must decide whether to initiate communication with its neighbors. Unbounded gossip would saturate the low-power wireless links that characterize edge deployments; conversely, suppressing all communication would prevent the cooperative resource allocation that distinguishes our approach from purely local reactive control.

We resolve this tension with a \textit{rate-gated broadcast condition}. Node~$i$ transmits a \texttt{NegotiationMessage} carrying its urgency~$u_i$, current temperature reading~$T_i$, and estimated rate of change~$\hat{\dot{T}}_i$ if and only if:
\begin{equation}
    |\hat{\dot{T}}_i| > \tau_{\text{talk}},
    \label{eq:talk_threshold}
\end{equation}
where $\tau_{\text{talk}} = 0.01\,^\circ\text{C/s}$ is the configurable talk threshold. This condition has two desirable properties. First, it suppresses communication entirely during steady-state periods when temperatures are stable, reducing the gossip message count to zero and conserving bandwidth. Second, it ensures that nodes experiencing rapid thermal transients---the precise situations in which coordinated resource reallocation is most valuable---reliably announce their state to the mesh. The threshold is intentionally conservative: at $0.01\,^\circ\text{C/s}$, a room would breach a $2\,^\circ\text{C}$ comfort margin in approximately 200~seconds, providing ample lead time for the gossip to propagate and the allocation to converge before a comfort violation occurs.

Each \texttt{NegotiationMessage} is a lightweight data structure comprising the origin node identifier, the origin's grid position, the relaying sender's identifier, a timestamp, the current scalar reading, the rate of change, the computed urgency, and a hop counter. This domain-agnostic design decouples the gossip layer from the physical quantity being monitored---the same protocol could carry CO\textsubscript{2} concentration urgencies, humidity urgencies, or power-demand urgencies without modification.

\subsection{Multi-Hop Propagation}
The gossip protocol operates over a two-layer architecture comprising a \textit{topology layer} (\texttt{SensorNetwork}) and an \textit{orchestration layer} (\texttt{SensorField}). \\
\subsubsection{\textbf{Topology Layer}} 
The \texttt{SensorNetwork} module constructs a communication graph $G = (V, E)$ using NetworkX, where each vertex $v \in V$ corresponds to a physical sensor node placed on the building grid. Sensor placement follows a configurable strategy (uniform grid with spacing~$s$); for the university floor plan, $s = 4\,\text{m}$ with sensors deployed at $z = 1$ (ceiling height), yielding $|V| = 64$ nodes across the five thermal zones. An undirected edge $(v_i, v_j) \in E$ is added whenever the Euclidean distance satisfies:
\begin{equation}
    \| \mathbf{p}_i - \mathbf{p}_j \| \cdot \Delta x \leq r_{\text{comm}} \cdot \Delta x,
    \label{eq:comm_radius}
\end{equation}
where $r_{\text{comm}} = 6\,\text{m}$ is the communication radius and $\Delta x = 1\,\text{m}$ is the grid cell size. The resulting graph exposes standard topology metrics---degree distribution, clustering coefficient, diameter, and average path length---used for academic reporting.
\subsubsection{\textbf{Orchestration Layer}}
The \texttt{SensorField} module orchestrates the three-phase \textit{sense $\rightarrow$ predict $\rightarrow$ gossip} pipeline on each simulation tick. A single call to \texttt{SensorField.step()} executes the following sequence:
\begin{enumerate}
    \item \textbf{Sense.} Each node's inbox is cleared, and new scalar readings $\{(v_i, T_i)\}$ are ingested via \texttt{node.sense()}, updating the rolling buffer and computing the least-squares slope $\hat{\dot{T}}_i$.
    \item \textbf{Predict.} Each node's TTI and urgency properties are implicitly updated as derived quantities of the refreshed buffer state.
    \item \textbf{Gossip.} The \texttt{\_run\_gossip()} method executes $R_{\text{gossip}}$ rounds of message propagation over the communication graph.
\end{enumerate}
During gossip initiation, each node~$i$ satisfying the broadcast condition~(\ref{eq:talk_threshold}) creates a \texttt{NegotiationMessage} and enqueues it for delivery to all graph neighbors $\mathcal{N}(i)$ returned by \texttt{SensorNetwork.neighbors()}. In each subsequent gossip round, the orchestrator iterates over all pending messages. When node~$j$ receives a message originating from node~$i$, it compares the incoming urgency against its stored value for that origin:
\begin{equation}
    u_i^{\text{new}} > u_i^{\text{stored}} \implies \text{accept, forward to } \mathcal{N}(j) \setminus \{\text{sender}\}.
    \label{eq:selective_forward}
\end{equation}
This selective forwarding ensures that only monotonically increasing urgency updates propagate through the network, preventing message storms while guaranteeing convergence. Forwarded messages carry an incremented hop counter, and messages exceeding $h_{\max} = 10$ hops are discarded as a safety bound, though in practice convergence occurs well within~$R_{\text{gossip}} = 2$ rounds.
After two gossip rounds with $r_{\text{comm}} = 6\,\text{m}$ and $s = 4\,\text{m}$ spacing, each node has received urgency estimates from all neighbors within two hops---sufficient to propagate information across any individual thermal zone and into adjacent zones through hallway connections. The per-zone urgency used for damper allocation is the maximum received urgency within the zone, ensuring that the most thermally stressed sensor within a large open-plan space drives the cooling allocation for that entire zone.

\subsection{Proportional Damper Allocation (Edge Policy)}

The shared cooling plant capacity~$Q_{\text{total}}$ is distributed proportionally to per-zone urgency. For zone~$i$ with local urgency~$u_i$, the allocation fraction and cooling assignment are:
\begin{equation}
    A_i = \frac{u_i}{\displaystyle\sum_{j} u_j}, \qquad Q_{\text{cool},i} = A_i \cdot Q_{\text{total}}.
    \label{eq:proportional_allocation}
\end{equation}

This formula transfers cooling capacity from low-urgency zones (empty rooms,
thermally stable corridors) to high-urgency zones (actively warming occupied
spaces), routing the contentioned resource budget to wherever it is most
imminently needed.  Unlike push-sum or averaging gossip protocols, which seek
global consensus on a scalar aggregate, our gossip propagates urgency as a
\emph{maximum-filter}: each node stores the highest urgency it has ever heard
from each origin, forwarding a message only when the incoming urgency
exceeds its stored value (see~\eqref{eq:selective_forward}).  This ensures
monotonically non-decreasing urgency estimates at each node and prevents
redundant message floods, at the cost that urgencies decay only through the
natural dropout of gossip once a zone's $|\hat{\dot{T}}|$ falls below
$\tau_{\text{talk}}$.

A critical implementation detail prevents pathological overcooling in the edge policy. Because each node's urgency is computed from local sensor readings with zero delay ($\text{AoI} = 0$), a zone with even a minuscule positive~$\hat{\dot{T}}$ could receive the entire cooling budget if all other zones report zero urgency. To address this, the damper opening is additionally capped by the absolute urgency:
\begin{equation}
    A_i^{\text{capped}} = \min\!\left(\frac{u_i}{\displaystyle\sum_{j} u_j},\; \kappa \cdot u_i\right),
    \label{eq:urgency_cap}
\end{equation}
where $\kappa = 20$ is a scaling constant selected empirically such that a zone with $\text{TTI} \leq 50$~seconds can access 100\% of its proportional share, while zones with $\text{TTI} > 100$~seconds are progressively throttled. This cap prevents low-urgency zones from monopolizing the cooling plant merely because their urgency happens to dominate in relative terms.

\subsection{Centralized Baseline}

To isolate the causal effect of data freshness on control performance, we implement a centralized baseline that uses the \textit{identical} proportional allocation formula~(\ref{eq:proportional_allocation}). A simulated Building Management System~(BMS) controller polls all zone sensors every $t_{\text{poll}}$~seconds. Each poll cycle introduces Gaussian network jitter $\mathcal{N}(0, \sigma_{\text{jitter}}^2)$ on the polling interval and a fixed supervisory compute delay~$d_{\text{compute}}$, yielding an effective Age of Information:
\begin{equation}
    \text{AoI}_{\text{cent}} = t_{\text{poll}} + \mathcal{N}(0, \sigma_{\text{jitter}}^2) + d_{\text{compute}}.
    \label{eq:aoi_centralized}
\end{equation}

Between polling events, the centralized controller issues actuation commands based on cached (stale) urgency values; damper positions are held constant until the next successful poll completes. This mirrors the behavior of real BACnet/IP polling architectures in commercial buildings, where supervisory controllers typically operate on 5--60~second polling cycles with non-negligible round-trip latencies~\cite{bacnet2016}.

Because both the edge and centralized policies use the identical allocation formula, any difference in comfort violation, energy consumption, or overshoot between the two modes can be attributed exclusively to the difference in data freshness---$\text{AoI} = 0$ for the edge policy versus $\text{AoI} \approx t_{\text{poll}} + d_{\text{compute}}$ for the centralized baseline. This controlled experimental design enables a clean causal analysis of the relationship between Age of Information and HVAC control quality.

\subsection{Evaluation Metrics}
\begin{table}[t]
\centering
\caption{Evaluation Metrics}
\label{tab:metrics}
\begin{tabular}{lll}
\toprule
\textbf{Metric} & \textbf{Definition} & \textbf{Unit} \\
\midrule
Comfort Violation & $\int \sum_i \max(T_i - T_{{\rm sp},i},\, 0)\, dt$ & °C·s \\
Cumulative Energy & $\int Q_{\rm cool}\, dt$ & W·s/m$^3$ \\
Max Overshoot     & $\max_{t,i}(T_i(t) - T_{{\rm sp},i})$ & °C \\
Mean AoI          & Average data age at actuation & s \\
Total Messages    & Cumulative gossip transmissions & --- \\
\bottomrule
\end{tabular}
\end{table}

Table~\ref{tab:metrics} summarizes the metrics used to compare the edge and centralized policies across all benchmark runs.
Comfort violation and energy are accumulated as running integrals over all
steps.  AoI is recorded per actuation event and averaged over the run.
Message count aggregates all gossip broadcasts; both policies run the same
gossip protocol on every timestep (the centralized baseline does not generate
poll messages in simulation --- it reads urgency from in-memory sensor state
at poll events).  Both policies run from identical initial conditions (same
seed, same environment configuration, same initial temperature field) to
ensure a controlled comparison.

\subsection{Simulation Modes}

\subsubsection{\textbf{Topology Layer}} 

Both policies share the identical top-level loop. Each tick of \texttt{Simulation.step()} executes two operations: (1)~\texttt{world.step()} advances the thermal physics by~$\Delta t$ (diffusion, heat injection, cooling), and (2)~\texttt{interface.step(world)} runs the full \textit{read\,$\rightarrow$\,infer\,$\rightarrow$\,actuate} cycle. The interface is the only component bridging the physical \texttt{World} and the inference \texttt{SensorField}---sensor nodes never observe the temperature field directly.

Within \texttt{interface.step()}, three sub-steps execute sequentially:
\begin{enumerate}
    \item \textbf{Read.} Extract scalar temperatures $T_i = \texttt{world.T}[\mathbf{p}_i]$ at each sensor position.
    \item \textbf{Infer.} Feed readings to \texttt{sensor\_field.step()}, which updates rolling buffers, recomputes slopes~$\hat{\dot{T}}_i$ and urgencies, then runs gossip for $R_{\text{gossip}}=2$ rounds over the communication graph.
    \item \textbf{Actuate.} Dispatch to either \texttt{\_evaluate\_edge()} or \texttt{\_evaluate\_centralized()} and apply the resulting damper openings via \texttt{world.set\_damper()}.
\end{enumerate}

Steps~1 and~2 are identical for both policies. The \textit{only} divergence occurs in Step~3, isolating data freshness as the sole experimental variable.
For the Edge mode,\texttt{\_evaluate\_edge()} computes damper openings using the two-stage urgency estimation (Section~IV-C): local urgency from co-located sensors, global denominator from gossip. Because the urgency values were computed from readings taken \textit{in the same tick}, $\text{AoI} = 0$. Each damper independently adjusts its opening every~$\Delta t$. For the Centralized mode, \texttt{\_evaluate\_centralized()} models a BMS polling architecture. A poll timer fires every $t_{\text{poll}} + \mathcal{N}(0, \sigma_{\text{jitter}}^2)$ seconds; between polls, the controller returns \textit{cached} openings from the previous cycle, holding all dampers fixed regardless of thermal transients. When a poll occurs, the controller reads sensor urgencies directly (ignoring gossip) and computes the same proportional allocation, recording $\text{AoI} = (t - t_{\text{last poll}}) + d_{\text{compute}}$. A sudden occupancy spike immediately after a poll will not be addressed for up to $t_{\text{poll}} + d_{\text{compute}}$ seconds --- the fundamental latency disadvantage that the edge architecture eliminates.



% ============================================================================
\section{Experimental Results}
% ============================================================================

\subsection{Experimental Setup}

All experiments use the \texttt{university\_floor.json} environment: five
thermal zones on a $46 \times 30 \times 3$ voxel grid at 1\,m resolution,
with thermal diffusivity $\alpha = 0.02\,\text{m}^2/\text{s}$ and an ambient
temperature of $20\,^\circ\text{C}$.  Zone setpoints are
$22\,^\circ\text{C}$ (classrooms, conference room, faculty office) and
$20\,^\circ\text{C}$ (computer lab).  Each zone is equipped with one VAV
damper; per-zone maximum flow capacities sum to $\sum_i f_i^{\max} = 7.5$,
and the shared plant capacity is set to $Q_{\text{total}} = 4.5$, enforcing
genuine resource contention ($Q_{\text{total}} / \sum f_i^{\max} \approx
60\%$).  Sensors are placed on a $4\,\text{m}$ grid at ceiling height
($z = 1$) with a communication radius of $r_{\text{comm}} = 6.0\,\text{m}$,
yielding \todo{N} nodes in a connected mesh.  All runs use
$\text{seed} = 42$ and 400 simulation steps.

Centralized baseline parameters: $t_{\text{poll}} = 15\,\text{s}$,
$\sigma_{\text{jitter}} = 2\,\text{s}$, $d_{\text{compute}} = 3\,\text{s}$,
producing a mean AoI of \todo{RERUN}\,s.  Edge parameters:
\texttt{gossip\_rounds} = 2, $\tau_{\text{talk}} = 0.01\,^\circ\text{C/s}$,
buffer window $= 30\,\text{s}$.

\subsection{Main Results}

Table~\ref{tab:results} summarizes the benchmark comparison.  Aether-Edge
outperforms the centralized baseline on four of five metrics, with the
performance gap attributed directly to the AoI difference.

\begin{table}[t]
\centering
\caption{Aether-Edge vs.\ Centralized Baseline
         (university\_floor, 400 steps, seed\,=\,42) --- \emph{pending re-run}}
\label{tab:results}
\begin{tabular}{lrrl}
\toprule
\textbf{Metric} & \textbf{Edge} & \textbf{Cent.} & \textbf{Gain} \\
\midrule
Comfort Violation  & \todo{--} & \todo{--} & \todo{--} \\
Energy             & \todo{--} & \todo{--} & \todo{--} \\
Max Overshoot (°C) & \todo{--} & \todo{--} & \todo{--} \\
Mean AoI (s)       & 0.00      & \todo{--} & \textbf{100\%} \\
Total Messages     & \todo{--} & \todo{--} & \todo{--} \\
\bottomrule
\end{tabular}
\end{table}

The \todo{X}\% reduction in comfort violation is the central result.  The
edge policy's AoI\,=\,0 means that urgency signals are computed from
readings taken \emph{in the same timestep} as the actuation decision.
Combined with TTI prediction — which generates a non-zero urgency signal
\emph{before} the setpoint is breached, proportional to how fast the zone
is warming — the edge policy pre-allocates cooling capacity to a rising zone
while the centralized baseline's dampers remain frozen at their
last-polled positions.  With $t_{\text{poll}} = 15\,\text{s}$,
$d_{\text{compute}} = 3\,\text{s}$, the centralized controller's actuation
lags the thermal event by up to \todo{W}\,s mean AoI.

The \todo{Z}\% difference in total message count between the two policies
reflects a difference in thermal state quality, not architectural
communication style.  Both policies execute the identical gossip protocol
on every timestep.  The centralized world's weaker thermal control produces
more frequent and larger temperature transients, causing more sensor nodes
to exceed $\tau_{\text{talk}}$ and generate gossip traffic.  The message
reduction is therefore a \emph{consequence} of better control, not a
separate design choice: a well-controlled environment with stable
temperatures is naturally a quiet network.

Energy consumption is \todo{sign and magnitude} between the two policies
(\todo{Y}\% difference), which is expected: both policies use the same
proportional allocation formula~\eqref{eq:proportional_allocation}, and the
shared cooling budget is fully consumed each step regardless of allocation.
Any difference arises from the overshoot-and-correct cycles that the
centralized policy incurs: when a poll-cycle lag allows a zone to overshoot,
the subsequent aggressive correction consumes energy that would not have been
needed under pre-emptive edge actuation.

The contention design ($Q_{\text{total}} = 4.5 < \sum f_i^{\max} = 7.5$)
means not all zones can receive full cooling simultaneously; the urgency
allocation must actively divert the scarce budget toward the most
thermally urgent zones at the expense of others.

\subsection{Time-Series Analysis}

The cumulative comfort violation curve for the centralized policy grows
steeply during the initial occupancy ramp in Classroom~101 (morning lecture
schedule).  The centralized controller's first poll fires after
$t_{\text{poll}} = 15\,\text{s}$, but the urgency computed from that stale
snapshot underestimates the zone's current warming rate, producing an
insufficient cooling command.  The edge policy detects the same ramp within
one buffer window ($t_{\text{buf}} = 30\,\text{s}$) and increases the
Classroom~101 damper opening in the next timestep, limiting early
accumulation.

Beyond the initial ramp, both curves continue to grow, but the centralized
policy's slope is consistently steeper: each $15\,\text{s}$ poll-cycle lag
allows a recurring unchecked warm period that accumulates over 400 steps
into the \todo{X}-fold total violation gap visible in
Table~\ref{tab:results}.

The energy curves for both policies grow roughly linearly and remain
\todo{close / diverging} throughout the simulation.
\todo{Update time-series analysis after benchmark re-run.}

\subsection{Sensitivity to Communication Parameters}

A planned parameter sweep over centralized polling interval
($t_{\text{poll}} \in \{2, 5, 10, 20\}$\,s), gossip rounds
($\{1, 2, 3, 5\}$), and sensor spacing ($\{2, 3, 5, 8\}$\,m) has been
designed using the \texttt{ExperimentRunner} infrastructure in
\texttt{configs/experiments/sweep\_sensor\_density.json}.  Based on the AoI
analysis, we predict that:
\begin{itemize}[]
  \item Comfort violation will scale proportionally with polling interval in
        the centralized baseline, as the mean AoI grows linearly with
        $t_{\text{poll}}$.
  \item Edge performance will degrade gracefully with sparser sensor
        spacing, as noisier TTI estimates (fewer samples per zone) widen
        the effective prediction horizon uncertainty.
  \item Two gossip rounds are sufficient for the university floor topology
        (diameter $\approx 4$ hops); larger buildings will require more rounds.
  \item The AoI\,=\,0 advantage of the edge policy is topology-independent
        and will hold at all polling intervals tested.
\end{itemize}

\subsection{Qualitative Heatmap Analysis}

Frame-by-frame inspection of the simulation heatmaps reveals a qualitatively
distinct actuation pattern between the two policies.  Under the edge policy,
damper openings adjust smoothly and continuously: within 1--2 simulation steps
of a zone's temperature beginning to rise, the corresponding damper opens in
proportion to the computed urgency, and the temperature field remains
spatially uniform near the setpoint throughout the run.

Under the centralized policy, visible warm spots develop and persist for
several frames between poll events.  Damper corrections arrive in discrete
jumps at each poll instant rather than as smooth responses, and the overshoot
visible between corrections contributes directly to the accumulated comfort
violation.  In the Computer Lab zone, which has a continuous baseline heat
load, the centralized controller's stale data consistently underestimates the
zone's current temperature during the quiet periods between polls, leading to
systematic under-cooling.

% ============================================================================
\section{Conclusion}
% ============================================================================

\subsection{Summary}

We built and evaluated Aether-Edge, a simulation framework demonstrating
urgency-based $n$-hop gossip for contentioned shared resource allocation.
The framework comprises a 3D FTCS thermal physics engine (used as the
demonstration environment), a per-node TTI inference engine based on OLS
rolling-buffer regression, an event-triggered multi-hop urgency gossip
protocol, a proportional resource allocator, a centralized polling baseline,
and a benchmark infrastructure with 190 unit and integration tests.

The core contribution is the urgency-gossip allocation algorithm: each node
independently computes urgency from local rate-of-change monitoring, gossips
it $n$ hops through a mesh topology, and allocates its proportional share of
a contentioned shared resource with AoI\,=\,0.  The demonstration on HVAC
shows that eliminating Age of Information --- from \todo{W}\,s (centralized)
to 0\,s (edge) --- produces a \todo{X}\% reduction in cumulative comfort
violation and \todo{Z}\% fewer gossip messages under genuine resource
contention ($Q_{\text{total}} < \sum f_i^{\max}$).  \todo{Update with
re-run results.}

\subsection{Limitations}

Several scope boundaries should be noted.  The thermal model solves linear
diffusion only; it does not capture advection, radiation exchange, humidity,
thermal mass, or nonlinear airflow effects.  The system is
simulation-only with no integration to BACnet, Modbus, or real sensor
hardware.  The two-round gossip approximation does not guarantee optimal
global consensus in sparse or disconnected graphs.  Occupancy schedules are
fixed and known; the system does not respond to real-time occupancy changes
beyond what the rolling buffer detects in the temperature signal.  Finally,
the building modeled has a single shared cooling plant; real installations
with per-zone Air Handling Units have additional capacity constraints not
modeled here.

\subsection{Future Work}

Six concrete extensions are prioritized for future work:
\begin{enumerate}[]
  \item \textbf{Hardware deployment} on Raspberry Pi Zero W nodes with DS18B20
        temperature sensors and relay-controlled VAV dampers, to validate the
        TTI inference latency on embedded hardware.
  \item \textbf{High-fidelity physics} via EnergyPlus or TRNSYS co-simulation
        through an FMI interface, replacing the FTCS engine with a calibrated
        whole-building model.
  \item \textbf{Occupancy-aware prediction} by integrating CO$_2$ or PIR
        sensor streams to update $Q_{\text{heat}}$ estimates in real time,
        extending the prediction horizon beyond what temperature rate-of-change
        alone provides.
  \item \textbf{AoI-optimal talk threshold} scheduling per Kaul and Yates
        \cite{kaul2020}, replacing the fixed $\tau_{\text{talk}}$ with a
        dynamically adapted threshold that minimizes AoI per unit transmitted
        energy.
  \item \textbf{Hierarchical two-tier gossip} for multi-plant buildings, where
        zone-level gossip propagates urgency within a floor and a floor-level
        gossip layer negotiates capacity allocation across plants.
  \item \textbf{Reinforcement learning baseline comparison} using a Soft
        Actor-Critic (SAC) agent trained in the same simulation environment, to
        quantify the gap between TTI heuristics and learned optimal control.
\end{enumerate}

% ============================================================================
% References
% ============================================================================

\begin{thebibliography}{99}

\bibitem{iea2023}
International Energy Agency, ``Buildings,'' IEA, Paris, 2023. [Online].
Available: https://www.iea.org/energy-system/buildings

\bibitem{ashrae2017}
ASHRAE, ``Standard 55-2017: Thermal Environmental Conditions for Human
Occupancy,'' American Society of Heating, Refrigerating and Air-Conditioning
Engineers, Atlanta, GA, 2017.

\bibitem{kaul2012}
S.~Kaul, R.~Yates, and M.~Gruteser, ``Real-time status: How often should one
update?,'' in \textit{Proc. IEEE INFOCOM}, Orlando, FL, Mar.\ 2012,
pp.~2731--2735.

\bibitem{oldewurtel2012}
F.~Oldewurtel, A.~Parisio, C.~N.~Jones, D.~Gyalistras, M.~Gwerder,
V.~Stauch, B.~Lehmann, and M.~Morari, ``Use of model predictive control and
weather forecasts for energy efficient building climate control,''
\textit{Energy and Buildings}, vol.~45, pp.~15--27, 2012.

\bibitem{ma2012}
Y.~Ma, F.~Borrelli, B.~Hencey, B.~Coffey, S.~Bengea, and P.~Haves,
``Model predictive control for the operation of building cooling systems,''
\textit{IEEE Trans.\ Control Syst.\ Technol.}, vol.~20, no.~3,
pp.~796--803, May 2012.

\bibitem{hao2012}
H.~Hao, B.~M.~Sanandaji, K.~Poolla, and T.~L.~Vincent, ``Aggregate
flexibility of thermostatically controlled loads,'' \textit{IEEE Trans.\
Power Syst.}, vol.~30, no.~1, pp.~189--198, Jan.\ 2015.

\bibitem{olfati2007}
R.~Olfati-Saber, J.~A.~Fax, and R.~M.~Murray, ``Consensus and cooperation in
networked multi-agent systems,'' \textit{Proc.\ IEEE}, vol.~95, no.~1,
pp.~215--233, Jan.\ 2007.

\bibitem{sun2017}
Y.~Sun, Y.~Polyanskiy, and E.~Uysal-Biyikoglu, ``Update or wait: How to keep
your information fresh,'' in \textit{Proc.\ IEEE INFOCOM}, Atlanta, GA,
May 2017.

\bibitem{kaul2020}
S.~Kaul and R.~D.~Yates, ``Age of information: Updates with priority,''
in \textit{Proc.\ IEEE INFOCOM Workshops}, Toronto, ON, Canada, Jul.\ 2020.

\bibitem{costa2016}
M.~Costa and A.~Ephremides, ``Age of information in wireless sensor networks
with channel randomness,'' in \textit{Proc.\ IEEE MILCOM}, Baltimore, MD,
Nov.\ 2016.

\bibitem{kempe2003}
D.~Kempe, A.~Dobra, and J.~Gehrke, ``Gossip-based computation of aggregate
information,'' in \textit{Proc.\ 44th Annu.\ IEEE Symp.\ Found.\ Comput.\
Sci.\ (FOCS)}, Cambridge, MA, 2003, pp.~482--491.

\bibitem{boyd2006}
S.~Boyd, A.~Ghosh, B.~Prabhakar, and D.~Shah, ``Randomized gossip
algorithms,'' \textit{IEEE Trans.\ Inf.\ Theory}, vol.~52, no.~6,
pp.~2508--2530, Jun.\ 2006.

\bibitem{rabbat2004}
M.~Rabbat and R.~Nowak, ``Distributed optimization in sensor networks,''
in \textit{Proc.\ 3rd Int.\ Symp.\ Inf.\ Process.\ Sensor Netw.\ (IPSN)},
Berkeley, CA, Apr.\ 2004, pp.~20--27.

\end{thebibliography}

\end{document}